% !TEX root = ../Abschlussarbeit_EN.tex
\section{Introduction}

For many people getting quickly and comfortably from A to B is a huge part of
their daily schedules.
Given the varied ways of transportation either by bike, bus, tram or car as
well as the underlying regulations and infrastructure a complex and multi-layered
traffic system is created.
To ensure effective and functional traffic management extensive data about a
given area is essential for numerous applications for example route finding,
real-time navigation or availability of different paths.
This is particularly relevant when looking at new technologies such as autonomous
driving.
Here, the vehicle needs great amounts of detailed data to successfully analyse
its environment so that safe and reliable travel can be guaranteed.
Remote sensing methods, especially the analysis of satellite imagery, are often
used in this context, to get valuable information about traffic infrastructure
and its surrounding area. Using the aerial view, many environmental features can
be detected, for example structures like roundabouts or intersections. Although
many attributes can be examined, this perspective also has its limits,
especially when looking at smaller traffic features, such as street signs or
road inclinations. To overcome this limitation, other technologies are used to
get information about the surrounding area from different perspectives.
Examples of that are radar-based concepts like LIDAR or images at ground level,
like street view images, which provide 360 degree images at a given location.

With remote sensing and complementary technologies providing increasingly
complex data for traffic management, the tools to analyse this information have
undergone substantial developments in recent years.
Since the launch of \textit{OpenAI’s ChatGPT} at the end of 2022, many
industries as well as research areas have seen rapid, ongoing developments when
using artificial intelligence (AI).
The technology can be found in almost every domain, for example in chatbots, AI
assistants or as a support tool in analysing different types of data.
The integration of different tools and models can be seen not only in IT, but
also in medicine, administration as well as transportation. For example, the
\textit{Fraunhofer IOSB-INA} is conducting research in traffic flow optimization using
AI technologies, as part of the \textit{AI4LSA} project \cite{FIOSBINA}.

In this project, we aim to combine environmental traffic data from different
sensing perspectives with artificial intelligence and deep learning methods,
to identify one-way streets at a given location. To determine the type of a
given street, the analysis is based on the detection of street signs, as they
are one of the primary indicators of one-way streets.
For that, the main approach is to analyse satellite images and street view
panoramas at intersections, using different AI models and techniques.


